\documentclass[10pt]{article}
\usepackage{icmlko}

\icmlmeta{IP 파운데이션 월드모델}{216}{관문이 잡음을 봉우리로 읽었다}{챔피언이 된 F21 을 열었더니 손잡이 하나가 죽어 있었다. 그것을 통과시킨 관문을 고치려다 두 번 실패하고, 되는 것 하나를 찾았다}

\begin{document}

\twocolumn[
\icmlheader

\begin{abstract}
노트 215가 챔피언을 F18 배깅에서 F21 로 넘기고 ``F21 은 한 번도 분해된
적이 없다''를 남겼다. 열었다. \textbf{손잡이 둘 중 하나가 죽어 있다} ---
최근 가중이 $+$0.0013($t{=}1.04$)이고 도메인별 풀링 선택이
$+$0.0105($t{=}2.67$)인데, \textbf{둘 다 켜면 $+$0.0105 로 선택만 켠
것과 소수 넷째 자리까지 같다}(차 $=$ 0.0000 · $t{=}0.00$). 노트 201이
``둘이 거의 가법($+$0.020 · $+$0.013)''이라고 한 것이 현재 상태에서
\textbf{반증된다.} 그리고 F6 의 알파 조율도 죽었다 --- 조율판이 $+$0.4286
으로 맨몸 $+$0.4288 보다 \emph{낮다.} \textbf{그런데 둘 다 노트 190의
안쪽 관문을 통과한다}(봉우리 0.585 · 0.628). 까닭은 봉우리 지수가
``(최고 $-$ 중앙) $\div$ (최고 $-$ 최저)''라 \textbf{상대 척도}이기
때문이다 --- 곡면 전체가 절대적으로 평평해도(범위 0.0058 · 0.0062) 1 에
가까운 값이 나온다. 고치려고 두 번 시도했고 두 번 다 실패했다.
① \textbf{주변 잡음 바닥} --- 같은 격자점을 다시 재면 sd 가 0.0132
$\sim$0.0136 이라 곡면이 통째로 잡음 안이지만, 격자점들은 \emph{같은}
안쪽 시험 표본에서 재므로 짝지어져 있어 주변 잡음을 쓰면 살아 있는
손잡이까지 막는다. ② \textbf{짝 잡음 바닥} --- 짝 sd 는 0.0028$\sim$0.0035
로 옳은 눈금인데, \textbf{그 $t$ 가 분할마다 1.77 과 3.07 을 오간다} ---
같은 안쪽 분할의 잡음을 물려받는다. 판도 안 움직였다. 되는 것은 셋째다:
\textbf{분할을 바꿔 보는 것.} 다섯 분할에서 최고가 알파는
{[}1, 100, 500, 500, 2{]}, $\tau$ 는 {[}없음, 8, 12, 1, 12{]} 로 \textbf{격자
전체를 헤맨다} --- 둘 다 최빈이 \textbf{2/5} 다. 봉우리 지수 자체도
0.035 에서 0.966 까지 흔들린다. \textbf{``고를 것이 없다''를 유보를 한 번도
안 보고 말할 수 있다.} 채택했다 --- 다섯 분할 중 셋에서 같은 칸이 이겨야
고른다. \textbf{F6 조율이 $+$0.4286 $\to$ $+$0.4288 로 맨몸과 같아지고}
(알파 선택이 막힌다) F21 $+$0.4395 · F18 $+$0.4396 은 그대로다.
\textbf{죽은 손잡이 둘이 판에서 사라졌고 판은 안 내려갔다.}
\end{abstract}
]

\begin{figure*}[t]
\centering
\includegraphics[width=\textwidth]{figs/stable.pdf}
\caption{왼쪽 --- 분할을 바꾸면 최고가 격자를 헤맨다. 가운데 --- 관문
자체가 흔들린다. 오른쪽 --- 안쪽은 통과하고 바깥은 0 이다.}
\end{figure*}

\section{F21 의 손잡이 하나가 죽어 있다}

\begin{table}[t]
\centering\tsize
\caption{맨몸 능형(F6 $+$0.4288)에서 얹어 가며. 짝 붓스트랩 500 회.}
\begin{tabular}{lrrr}
\toprule
& 판 $\rho$ & 맨몸 대비 & $t$ \\
\midrule
최근 가중만 & .4302 & $+$.0013 & 1.04 \\
\textbf{풀링 선택만} & \textbf{.4395} & \textbf{$+$.0105} & \textbf{2.67} \\
F21 (둘 다) & .4395 & $+$.0105 & 2.67 \\
\midrule
\multicolumn{2}{l}{둘 $-$ 선택만} & \textbf{0.0000} & \textbf{0.00} \\
둘 $-$ 최근만 & & $+$.0092 & 2.42 \\
\bottomrule
\end{tabular}
\end{table}

\claimx{\textbf{F21 의 이득은 통째로 풀링 선택이다.} 최근 가중을 얹어도
소수 넷째 자리까지 안 움직인다.}{노트 201이 ``각각 $+$0.020 · $+$0.013
이고 합쳐서 $+$0.038 로 거의 가법''이라고 했다 --- \textbf{현재 상태에서
반증된다.} 그 사이 노트 207 · 210 · 213 · 214가 자료를 넷 고쳤다.}

\claimx{\textbf{F6 의 알파 조율도 죽었다} --- 조율판이 $+$0.4286 으로
맨몸 $+$0.4288 보다 낮다.}{노트 190이 ``능형 셋이 봉우리 0.81$\sim$0.84
로 전부 $+$0.011$\sim$0.013 을 얻는다''고 한 그 손잡이다. \textbf{살아
있는 손잡이가 아니라 살아 있던 손잡이였다.}}

\section{봉우리 지수는 상대 척도다}

\begin{table}[t]
\centering\tsize
\caption{안쪽 곡면. 문턱은 봉우리 0.5 또는 단조 0.8.}
\begin{tabular}{lrrrr}
\toprule
& 범위 & 봉우리 & 관문 & 바깥 \\
\midrule
능형 $\alpha$ & .0062 & \textbf{.628} & \textbf{통과} & $-$.0002 \\
감쇠 $\tau$ & .0058 & \textbf{.585} & \textbf{통과} & $+$.0013 \\
\bottomrule
\end{tabular}
\end{table}

\claimx{\textbf{봉우리 지수는 ``(최고 $-$ 중앙) $\div$ (최고 $-$ 최저)''
라 상대 척도다} --- 곡면이 절대적으로 평평해도 1 에 가까운 값이
나온다.}{범위가 0.006 이면 격자 안에서 고를 것이 없는데 봉우리 지수는
그것을 못 본다. \textbf{노트 190이 세운 규약의 구멍이고 스물여섯 노트
동안 안 보였다.}}

\section{두 번 실패했다}

\begin{table}[t]
\centering\tsize
\caption{잡음 바닥을 어떻게 재나.}
\begin{tabular}{lrrl}
\toprule
& $\alpha$ & $\tau$ & \\
\midrule
격자 범위 & .0062 & .0058 & \\
\textbf{주변} sd & .0136 & .0132 & 너무 세다 \\
\textbf{짝} sd & .0035 & .0028 & 눈금은 옳다 \\
짝 $t$ & 1.77 & 2.05 & \\
\bottomrule
\end{tabular}
\end{table}

\claimx{\textbf{주변 잡음은 못 쓴다} --- 격자점을 \emph{같은} 안쪽 시험
표본에서 재므로 짝지어져 있고, 주변 sd 는 짝 sd 의 넉 배쯤이라 살아 있는
손잡이까지 막는다.}{짝 sd 가 옳은 눈금이다. 그런데 그것으로 바닥을 세워도
안 된다 --- \textbf{$t$ 가 분할을 바꾸면 1.77 과 3.07 을 오간다.} 같은
안쪽 분할에서 재므로 그 분할의 잡음을 그대로 물려받는다.}

\claimx{\textbf{그리고 판이 안 움직였다}(F6 조율 $+$0.4286 그대로).
바닥이 실제 접힘 안에서는 안 걸렸다.}{\textbf{노트 187의 제일 값싼 검사가
또 맞았다} --- 무언가를 바꿨는데 판이 넷째 자리까지 안 움직이면 그건
효과가 없는 것이 아니라 \emph{안 닿은} 것이다. 여기서는 닿았지만
문턱을 못 넘었다.}

\section{되는 것은 분할을 바꾸는 것}

\begin{table}[t]
\centering\tsize
\caption{안쪽 분할 지점을 다섯 번 옮기면.}
\begin{tabular}{lrrrr}
\toprule
분할 & $\alpha$ 최고 & 봉우리 & $\tau$ 최고 & 봉우리 \\
\midrule
.60 & 1 & .035 & 없음 & .230 \\
.65 & 100 & \textbf{.966} & 8 & .569 \\
.70 & \textbf{500} & .628 & \textbf{12} & .585 \\
.75 & \textbf{500} & .709 & 1 & \textbf{.804} \\
.80 & 2 & .111 & \textbf{12} & .590 \\
\midrule
최빈 & \textbf{2/5} & & \textbf{2/5} & \\
\bottomrule
\end{tabular}
\end{table}

\claimx{\textbf{최고가 격자 전체를 헤맨다} --- 알파는 1 에서 500 까지
(격자가 0.5$\sim$500 이다), $\tau$ 는 없음에서 12 까지. 둘 다 최빈이
2/5 다.}{\textbf{봉우리 지수 자체도 0.035 에서 0.966 까지 흔들린다} ---
관문이 곡면을 재는 것이 아니라 \emph{그 분할}을 재고 있었다.}

\claimx{\textbf{유보를 한 번도 안 본다.} 안쪽 분할만 다섯 번 옮긴다 ---
값싸고(격자 $\times$ 5회 적합) 라벨을 안 쓴다.}{규약으로 적는다:
\textbf{다섯 분할 중 셋에서 같은 칸이 이겨야 고른다.} 격자가 일곱이면
우연히 셋이 겹칠 확률이 4\% 안이다.}

\begin{table}[t]
\centering\tsize
\caption{채택 뒤의 판.}
\begin{tabular}{lrr}
\toprule
& 전 & 후 \\
\midrule
F6 능형(맨몸) & .4288 & .4288 \\
\textbf{F6 조율} & .4286 & \textbf{.4288} \\
F21 & .4395 & .4395 \\
F18 배깅 & .4396 & .4396 \\
\bottomrule
\end{tabular}
\end{table}

\claimx{\textbf{죽은 손잡이 둘이 사라졌고 판은 안 내려갔다.} F6 조율이
맨몸과 같아지고(알파 선택이 막힌다) 나머지는 그대로다.}{\textbf{점수를
올리는 고침이 아니라 없는 손잡이를 없다고 적는 고침이다} --- 노트 210의
규칙대로 점수가 아니라 측정의 유효성으로 정한다.}

\begin{table}[t]
\centering\tsize
\caption{남은 것.}
\begin{tabular}{ll}
\toprule
\textbf{다른 격자} & F18 의 깊이 · 반복 격자도 같은 병인가 \\
\textbf{풀링 선택} & F21 에 남은 하나 --- 그것도 안정한가 \\
셋/다섯 & 문턱을 격자 크기로 맞출 것 \\
\bottomrule
\end{tabular}
\end{table}

둘째 줄이 제일 값이 크다. \textbf{F21 의 이득이 통째로 풀링 선택 하나에
있다면 그 하나가 안정한지를 같은 잣대로 봐야 한다.} 이 노트가 잰 것은
$\tau$ 격자이고, 풀링 선택은 도메인마다 ``혼자가 나은가''를 안쪽 격차로
정하는 \emph{다른} 결정이다. \textbf{분할을 다섯 번 옮겼을 때 어느
도메인이 뽑히는지가 그대로인가} --- 그것이 흔들리면 챔피언의 유일한
손잡이가 잡음이고, 그러면 F21 도 F6 맨몸과 구분되지 않는다.

\end{document}
